
\documentclass[11pt]{article}

\usepackage[english]{babel}
\usepackage[letterpaper, top=2.5cm, bottom=2.5cm, left=2.5cm, right=2.5cm]{geometry}
\usepackage{amsmath}
\usepackage{graphicx}
\usepackage{booktabs}
\usepackage{longtable}
\usepackage{pdflscape}
\usepackage{caption}
\usepackage{setspace}
\usepackage[colorlinks=true, allcolors=blue]{hyperref}

\title{%
  \textbf{Replication of Koijen \& Yogo (2020):}\\[4pt]
  \large Exchange Rates and Asset Prices in a Global Demand System
}
\author{Replication Project}
\date{\today}

\begin{document}
\maketitle
\tableofcontents
\newpage

%% ============================================================
%%  SECTION 1 — PROJECT OVERVIEW (auto-generated)
%% ============================================================
\section{Replication Overview}
\label{sec:overview}

This report documents the replication of Table~1 (Market Values of Financial
Assets) and Table~2 (Top Ten Investors by Asset Class) from Koijen \& Yogo
(2020), ``Exchange Rates and Asset Prices in a Global Demand System,'' NBER
Working Paper~27342.  The paper develops a global demand system for financial
assets in which exchange rates and asset prices are jointly determined by the
portfolio choices of investors across 37 countries.  The key empirical
contribution is a comprehensive dataset of outstanding amounts and bilateral
cross-border holdings for short-term debt, long-term debt, and equity.

\subsection{Where the Replication Succeeded}

The API-based pipeline achieved close numerical agreement with the paper for
most OECD developed-market economies.  For the 2020 cross-check year,
long-term debt amounts for France, Germany, Italy, Spain, Canada, and Japan
matched the paper's ground-truth values within 5\%, and the short-term debt
rankings in Table~2 reproduced the correct ordering of the top investors.
Approximately 63\% of individual cell-level comparisons passed a 15\%
tolerance test when benchmarked against the paper's dataverse values.

\subsection{Where the Replication Faced Challenges}

Four structural limitations prevented full replication with freely available
data:
\begin{enumerate}
  \item \textbf{Equity amounts.}  The paper uses proprietary Datastream
    market-capitalisation data.  The best free substitute---OECD Table~720
    instrument F5 (equity and investment fund shares)---overstates market
    capitalisation by roughly 2--3$\times$ because it includes investment
    fund shares.  This affects all equity cells in Table~1.
  \item \textbf{Reserve shares.}  The paper uses IMF COFER data to allocate
    official foreign-exchange reserves by currency and issuer.  The publicly
    available CPIS S121 (central bank) reporter list is too sparse
    ($\approx$8 reporters for the United States) to construct reliable
    reserve shares.
  \item \textbf{BIS currency adjustment.}  The BIS Domestic Debt Securities
    Statistics (WS\_NA\_SEC\_DSS) reports all-currency totals.  A subtraction
    of the foreign-currency slice from WS\_DEBT\_SEC2\_PUB was applied for
    most countries; however, for Finland and Germany the OECD T720 national
    accounts data already reflect local-currency-only amounts, so the BIS
    adjustment was skipped for those two countries.
  \item \textbf{Data gaps.}  New Zealand debt is absent from all free sources;
    China reports only long-term general-government debt to BIS, missing
    short-term and non-government sectors entirely.
\end{enumerate}

\newpage

%% ============================================================
%%  SECTION 2 — DATA SOURCES
%% ============================================================
\section{Data Sources}
\label{sec:data}

\subsection{OECD Table 720 (DF\_T720R\_A)}
OECD Table~720 reports national financial accounts balance sheet data for
27 OECD member countries, sourced via the OECD SDMX REST API in CSV format.
Financial instrument F3 (debt securities) is split into short-term (S) and
long-term (L) maturities.  Instrument F5 (equity and investment fund shares)
provides an equity proxy, though it is known to overstate market capitalisation.
Values are reported in millions of local currency; they are converted to USD
billions using end-of-year exchange rates embedded in the API response.

\subsection{BIS Domestic Debt Securities (WS\_NA\_SEC\_DSS)}
The BIS Domestic Debt Securities Statistics cover 37 countries at quarterly
frequency, reported in USD billions.  This series is used as the primary
debt-outstanding source for non-OECD economies (China, India, Malaysia,
Philippines, Russia, South Africa, Thailand, Hong Kong, Singapore, Australia)
and as a fallback for OECD members not covered by T720.  VALUATION differs
by country: most report at nominal value (N), Australia at market value (M),
and Thailand at face value (F).

\subsection{BIS International Debt Securities (WS\_DEBT\_SEC2\_PUB)}
The BIS IDS foreign-currency slice is used to convert all-currency domestic
totals to local-currency-only amounts, following the Appendix~C methodology
of Koijen \& Yogo (2020).  Finland and Germany are excluded from this
correction because their OECD T720 submissions already reflect
local-currency amounts.

\subsection{IMF Portfolio Investment Position (PIP/CPIS)}
The IMF PIP bilateral parquet provides cross-border holdings between investor
and issuer country pairs for short-term debt, long-term debt, and equity.
These data underpin the foreign-ownership shares in Table~1 and the investor
rankings in Table~2.  Values are in USD billions; the \texttt{value} column
from the SDMX API is used directly (the \texttt{value\_usd} column applies
an erroneous double-scaling).

\subsection{World Bank WDI}
World Bank Development Indicators (series CM.MKT.LCAP.CD) provide
total stock-market capitalisation in USD for non-OECD economies where
OECD T720 equity data are unavailable.  These are the primary equity source
for India, China, Malaysia, Philippines, Thailand, South Africa, and Russia.

\newpage

%% ============================================================
%%  SECTION 3 — PAPER REPLICATION (Stata dataverse)
%% ============================================================
\section{Paper Replication: Original Dataverse Data}
\label{sec:orig}

Tables~\ref{tab:t1_orig} and \ref{tab:t2_orig} reproduce Table~1 and Table~2
from Koijen \& Yogo (2020) using the authors' dataverse Stata files for the
year 2020.  These serve as the ground-truth benchmark for evaluating the
API-based replication.

\subsection{Table 1: Market Values of Financial Assets (Paper Replication)}

\begin{landscape}
\small
\begin{longtable}{l r r r r r r r r r}
\caption{Market values of financial assets, 2020. Replication of Table~1 in Koijen \& Yogo (2020) using the authors' dataverse Stata files. Columns report outstanding amounts in billion USD together with the share held by domestic investors and the share attributable to official reserve holdings. The United States dominates all asset classes by absolute size; the domestic share is high across most markets, consistent with the home-bias puzzle the paper seeks to explain.} \label{tab:t1_orig} \\
\hline\hline
 & \multicolumn{3}{c}{\textit{Short-term debt}} & \multicolumn{3}{c}{\textit{Long-term debt}} & \multicolumn{3}{c}{\textit{Equity}} \\
\cmidrule(lr){2-4}\cmidrule(lr){5-7}\cmidrule(lr){8-10}
Country & Bn USD & Dom. & Res. & Bn USD & Dom. & Res. & Bn USD & Dom. & Res. \\
\hline
\endfirsthead
\caption*{Market values of financial assets, 2020. Replication of Table~1 in Koijen \& Yogo (2020) using the authors' dataverse Stata files. Columns report outstanding amounts in billion USD together with the share held by domestic investors and the share attributable to official reserve holdings. The United States dominates all asset classes by absolute size; the domestic share is high across most markets, consistent with the home-bias puzzle the paper seeks to explain. (continued)} \\
\hline\hline
 & \multicolumn{3}{c}{\textit{Short-term debt}} & \multicolumn{3}{c}{\textit{Long-term debt}} & \multicolumn{3}{c}{\textit{Equity}} \\
\cmidrule(lr){2-4}\cmidrule(lr){5-7}\cmidrule(lr){8-10}
Country & Bn USD & Dom. & Res. & Bn USD & Dom. & Res. & Bn USD & Dom. & Res. \\
\hline
\endhead
\hline
\multicolumn{10}{r}{\footnotesize\textit{Continued on next page}} \\
\endfoot
\hline\hline
\endlastfoot
\end{longtable}
\end{landscape}

\subsection{Table 2: Top Ten Investors by Asset Class (Paper Replication)}

\begin{table}[ht]
\centering
\caption{Top ten investors by asset class, 2020 (paper replication). Rankings are based on the IMF CPIS bilateral holdings in the authors' dataverse. The United States is the largest foreign investor in all three classes; Japan is a distant second in short-term and long-term debt, while the ranking for equity is more dispersed. \textit{Reserves} denotes aggregate official reserve holdings.}
\label{tab:t2_orig}
\small
\begin{tabular}{l r l r l r}
\hline\hline
\multicolumn{2}{c}{\textit{Short-term debt}} & \multicolumn{2}{c}{\textit{Long-term debt}} & \multicolumn{2}{c}{\textit{Equity}} \\
\cmidrule(lr){1-2}\cmidrule(lr){3-4}\cmidrule(lr){5-6}
Country & Bn USD & Country & Bn USD & Country & Bn USD \\
\hline
\hline\hline
\end{tabular}
\end{table}

\newpage

%% ============================================================
%%  SECTION 4 — FIGURES (paper replication)
%% ============================================================
\section{Figures: Paper Replication}
\label{sec:figures}

\subsection{Figure 1: Short-Term Debt and Interest-Rate Differentials}
\begin{figure}[ht]
\centering
\includegraphics[width=0.92\textwidth]{../_output/figure1_replicated.png}
\caption{\textbf{Figure 1 replication: short-term debt and short-term interest-rate differentials.} Each panel plots the foreign short-term debt share (x-axis) against the short-term interest-rate differential relative to the United States (y-axis) for the Euro area, Japan, Switzerland, and the United Kingdom. A positive slope indicates that higher foreign rates attract larger cross-border debt inflows, consistent with the paper's demand-system framework.}
\label{fig:fig1}
\end{figure}

\subsection{Figure 2: Long-Term Debt and Yield Differentials}
\begin{figure}[ht]
\centering
\includegraphics[width=0.92\textwidth]{../_output/figure2_replicated.png}
\caption{\textbf{Figure 2 replication: long-term debt and long-term yield differentials.} Each panel plots the foreign long-term debt share (x-axis) against the 5-year zero-coupon yield differential relative to the United States (y-axis), estimated via Nelson-Siegel from 10-year par yields. The positive slope is steeper than in Figure~1, reflecting greater yield sensitivity in long-duration bond demand.}
\label{fig:fig2}
\end{figure}

\newpage

%% ============================================================
%%  SECTION 5 — API-BASED REPLICATION (latest data: 2024)
%% ============================================================
\section{API-Based Replication: Latest Data (2024)}
\label{sec:latest}

Tables~\ref{tab:t1_2024} and \ref{tab:t2_2024} present the replication
tables for the most recent data vintage (2024) sourced entirely from
publicly available APIs.  Compared with the 2020 benchmark, outstanding
amounts have grown substantially: US long-term debt rose from approximately
\$41 trillion to \$51 trillion, and US equity from \$81 trillion to
\$116 trillion, reflecting the decade-long expansion of global capital markets.

\subsection{Table 1: Market Values of Financial Assets (2024)}

\begin{landscape}
\small
\begin{longtable}{l r r r r r r r r r}
\caption{Market values of financial assets, 2024 (latest available data). Produced from the most recent API vintage as of the pull date. Compared with 2020, total outstanding amounts have grown substantially in most markets, with US long-term debt rising from roughly \$41 trillion to \$51 trillion and US equity from \$81 trillion to \$116 trillion. The domestic share patterns are broadly stable over time, indicating persistent home-bias structures.} \label{tab:t1_2024} \\
\hline\hline
 & \multicolumn{3}{c}{\textit{Short-term debt}} & \multicolumn{3}{c}{\textit{Long-term debt}} & \multicolumn{3}{c}{\textit{Equity}} \\
\cmidrule(lr){2-4}\cmidrule(lr){5-7}\cmidrule(lr){8-10}
Country & Bn USD & Dom. & Res. & Bn USD & Dom. & Res. & Bn USD & Dom. & Res. \\
\hline
\endfirsthead
\caption*{Market values of financial assets, 2024 (latest available data). Produced from the most recent API vintage as of the pull date. Compared with 2020, total outstanding amounts have grown substantially in most markets, with US long-term debt rising from roughly \$41 trillion to \$51 trillion and US equity from \$81 trillion to \$116 trillion. The domestic share patterns are broadly stable over time, indicating persistent home-bias structures. (continued)} \\
\hline\hline
 & \multicolumn{3}{c}{\textit{Short-term debt}} & \multicolumn{3}{c}{\textit{Long-term debt}} & \multicolumn{3}{c}{\textit{Equity}} \\
\cmidrule(lr){2-4}\cmidrule(lr){5-7}\cmidrule(lr){8-10}
Country & Bn USD & Dom. & Res. & Bn USD & Dom. & Res. & Bn USD & Dom. & Res. \\
\hline
\endhead
\hline
\multicolumn{10}{r}{\footnotesize\textit{Continued on next page}} \\
\endfoot
\hline\hline
\endlastfoot
\multicolumn{10}{l}{\textit{Developed markets: North America}} \\
Canada & 479 & 0.68 & 0.00 & 2,249 & 0.61 & 0.02 & 8,321 & 0.86 & 0.00 \\
United States & 7,113 & 0.96 & 0.00 & 50,162 & 0.90 & 0.01 & 116,168 & 0.92 & 0.00 \\
\multicolumn{10}{l}{\textit{Developed markets: Europe}} \\
Austria & 28 & 0.64 & 0.02 & 545 & 0.54 & 0.03 & 1,002 & 0.94 & 0.00 \\
Belgium & 104 & 0.68 & 0.03 & 740 & 0.58 & 0.04 & 1,788 & 0.93 & 0.00 \\
Finland & 63 & 0.74 & 0.02 & 363 & 0.50 & 0.05 & 595 & 0.75 & 0.00 \\
France & 652 & 0.79 & 0.01 & 4,675 & 0.69 & 0.02 & 13,245 & 0.91 & 0.00 \\
Germany & 262 & 0.55 & 0.03 & 3,817 & 0.75 & 0.01 & 8,180 & 0.88 & 0.00 \\
Greece & 9 & 0.87 & 0.01 & 135 & 0.90 & 0.01 & 299 & 0.92 & 0.00 \\
Italy & 142 & 0.69 & 0.07 & 3,266 & 0.82 & 0.02 & 3,825 & 0.93 & 0.00 \\
Netherlands & 78 & 0.26 & 0.02 & 1,320 & 0.25 & 0.11 & 7,001 & 0.87 & 0.00 \\
Portugal & 22 & 0.81 & 0.01 & 289 & 0.74 & 0.03 & 677 & 0.96 & 0.00 \\
Spain & 117 & 0.75 & 0.03 & 1,909 & 0.66 & 0.03 & 4,051 & 0.93 & 0.00 \\
Denmark & 3 & 0.00 & 0.18 & 539 & 0.73 & 0.03 & 2,464 & 0.85 & 0.00 \\
Norway & 44 & 0.66 & 0.00 & 202 & 0.21 & 0.06 & 1,005 & 0.92 & 0.00 \\
Sweden & 119 & 0.58 & 0.05 & 383 & 0.48 & 0.09 & 3,054 & 0.90 & 0.00 \\
Switzerland & 131 & 0.93 & 0.00 & 586 & 0.77 & 0.01 & 3,286 & 0.71 & 0.00 \\
United Kingdom & 316 & 0.58 & 0.02 & 3,010 & 0.62 & 0.01 & 6,754 & 0.68 & 0.00 \\
\multicolumn{10}{l}{\textit{Developed markets: Pacific}} \\
Australia & 279 & 0.74 & 0.00 & 1,445 & 0.66 & 0.03 & 1,737 & 0.72 & 0.00 \\
Hong Kong & 289 & 0.89 & 0.01 & 0 & 1.00 & 1.00 & 4,550 & 0.84 & 0.00 \\
Japan & 960 & 0.86 & 0.03 & 8,583 & 0.94 & 0.00 & 10,581 & 0.84 & 0.00 \\
New Zealand & 0 & 1.00 & 1.00 & 76 & 0.30 & 0.04 & 93 & 0.63 & 0.00 \\
Singapore & 312 & 0.81 & 0.00 & 235 & 0.42 & 0.02 & 638 & 0.64 & 0.00 \\
\multicolumn{10}{l}{\textit{Emerging markets}} \\
Brazil & 168 & 0.90 & 0.01 & 2,490 & 0.96 & 0.00 & 659 & 0.74 & 0.00 \\
Colombia & 42 & 1.00 & 0.00 & 146 & 0.71 & 0.01 & 591 & 0.99 & 0.00 \\
Czech Republic & 15 & 0.28 & 0.00 & 148 & 0.79 & 0.00 & 503 & 0.99 & 0.00 \\
Hungary & 22 & 0.73 & 0.00 & 104 & 0.66 & 0.01 & 437 & 0.97 & 0.00 \\
Malaysia & 13 & 0.70 & 0.07 & 418 & 0.88 & 0.01 & 449 & 0.90 & 0.00 \\
Mexico & 115 & 0.96 & 0.03 & 692 & 0.78 & 0.00 & 1,825 & 0.95 & 0.00 \\
Philippines & 27 & 0.96 & 0.01 & 0 & 1.00 & 1.00 & 252 & 0.90 & 0.00 \\
Poland & 74 & 1.00 & 0.00 & 318 & 0.76 & 0.01 & 693 & 0.96 & 0.00 \\
Russia & 4 & 0.98 & 0.01 & 470 & 0.99 & 0.00 & \multicolumn{1}{c}{---} & \multicolumn{1}{c}{---} & \multicolumn{1}{c}{---} \\
South Africa & 40 & 0.95 & 0.00 & 336 & 0.88 & 0.00 & 986 & 0.91 & 0.00 \\
South Korea & 269 & 0.92 & 0.00 & 1,824 & 0.90 & 0.02 & 2,413 & 0.86 & 0.00 \\
Thailand & 79 & 0.98 & 0.01 & 414 & 0.95 & 0.00 & 520 & 0.89 & 0.00 \\
\end{longtable}
\end{landscape}

\subsection{Table 2: Top Ten Investors by Asset Class (2024)}

\begin{table}[ht]
\centering
\caption{Top ten investors by asset class, 2024 (latest available data). The United States retains its dominant position across all three asset classes in the latest vintage. China has risen in the equity ranking, reflecting continued growth in Chinese institutional asset management. Values are in billion USD.}
\label{tab:t2_2024}
\small
\begin{tabular}{l r l r l r}
\hline\hline
\multicolumn{2}{c}{\textit{Short-term debt}} & \multicolumn{2}{c}{\textit{Long-term debt}} & \multicolumn{2}{c}{\textit{Equity}} \\
\cmidrule(lr){1-2}\cmidrule(lr){3-4}\cmidrule(lr){5-6}
Country & Bn USD & Country & Bn USD & Country & Bn USD \\
\hline
United States & 7,220 & United States & 48,185 & United States & 114,220 \\
Japan & 881 & Japan & 10,437 & France & 12,829 \\
France & 775 & France & 5,269 & China & 11,592 \\
United Kingdom & 634 & United Kingdom & 4,806 & Japan & 10,279 \\
Hong Kong & 441 & Germany & 4,617 & Canada & 9,218 \\
Singapore & 380 & Italy & 3,456 & Germany & 8,243 \\
Canada & 355 & Canada & 3,267 & Netherlands & 7,009 \\
Australia & 297 & Brazil & 2,489 & United Kingdom & 6,460 \\
South Korea & 275 & South Korea & 2,125 & India & 4,352 \\
Switzerland & 195 & Spain & 1,741 & Hong Kong & 4,349 \\
\hline\hline
\end{tabular}
\end{table}

\newpage

%% ============================================================
%%  SECTION 6 — API CROSS-CHECK AT 2020
%% ============================================================
\section{API Cross-Check: 2020}
\label{sec:crosscheck}

Tables~\ref{tab:t1_2020} and \ref{tab:t2_2020} present the API-based
replication for 2020, the same year as the paper.  Direct comparison with
Tables~\ref{tab:t1_orig} and \ref{tab:t2_orig} reveals where the free-data
pipeline agrees and where it diverges.  Long-term debt amounts are broadly
consistent for most OECD economies; equity amounts are overstated due to
the OECD F5 proxy issue described in Section~\ref{sec:overview}.

\subsection{Table 1: Market Values of Financial Assets (2020, API)}

\begin{landscape}
\small
\begin{longtable}{l r r r r r r r r r}
\caption{Market values of financial assets, 2020 (API replication). Same format as Table~\ref{tab:t1_orig} but sourced entirely from publicly available APIs (OECD T720, BIS WS\_NA\_SEC\_DSS, IMF PIP, World Bank WDI). This cross-check year allows direct comparison with the paper values to assess replication accuracy. Amounts are broadly consistent with the paper but differ where API coverage diverges from proprietary data sources.} \label{tab:t1_2020} \\
\hline\hline
 & \multicolumn{3}{c}{\textit{Short-term debt}} & \multicolumn{3}{c}{\textit{Long-term debt}} & \multicolumn{3}{c}{\textit{Equity}} \\
\cmidrule(lr){2-4}\cmidrule(lr){5-7}\cmidrule(lr){8-10}
Country & Bn USD & Dom. & Res. & Bn USD & Dom. & Res. & Bn USD & Dom. & Res. \\
\hline
\endfirsthead
\caption*{Market values of financial assets, 2020 (API replication). Same format as Table~\ref{tab:t1_orig} but sourced entirely from publicly available APIs (OECD T720, BIS WS\_NA\_SEC\_DSS, IMF PIP, World Bank WDI). This cross-check year allows direct comparison with the paper values to assess replication accuracy. Amounts are broadly consistent with the paper but differ where API coverage diverges from proprietary data sources. (continued)} \\
\hline\hline
 & \multicolumn{3}{c}{\textit{Short-term debt}} & \multicolumn{3}{c}{\textit{Long-term debt}} & \multicolumn{3}{c}{\textit{Equity}} \\
\cmidrule(lr){2-4}\cmidrule(lr){5-7}\cmidrule(lr){8-10}
Country & Bn USD & Dom. & Res. & Bn USD & Dom. & Res. & Bn USD & Dom. & Res. \\
\hline
\endhead
\hline
\multicolumn{10}{r}{\footnotesize\textit{Continued on next page}} \\
\endfoot
\hline\hline
\endlastfoot
\multicolumn{10}{l}{\textit{Developed markets: North America}} \\
Canada & 530 & 0.76 & 0.00 & 2,484 & 0.69 & 0.01 & 6,599 & 0.87 & 0.00 \\
United States & 5,703 & 0.96 & 0.00 & 39,964 & 0.89 & 0.01 & 80,837 & 0.92 & 0.00 \\
\multicolumn{10}{l}{\textit{Developed markets: Europe}} \\
Austria & 18 & 0.71 & 0.06 & 607 & 0.55 & 0.03 & 937 & 0.93 & 0.00 \\
Belgium & 83 & 0.69 & 0.04 & 861 & 0.59 & 0.02 & 1,912 & 0.93 & 0.00 \\
Finland & 67 & 0.80 & 0.06 & 370 & 0.49 & 0.04 & 692 & 0.77 & 0.00 \\
France & 644 & 0.85 & 0.01 & 5,169 & 0.69 & 0.02 & 11,799 & 0.90 & 0.00 \\
Germany & 354 & 0.76 & 0.02 & 4,306 & 0.76 & 0.02 & 8,209 & 0.88 & 0.00 \\
Greece & 14 & 0.66 & 0.02 & 116 & 0.87 & 0.02 & 204 & 0.93 & 0.00 \\
Italy & 158 & 0.59 & 0.08 & 3,671 & 0.80 & 0.02 & 3,133 & 0.93 & 0.00 \\
Netherlands & 89 & 0.42 & 0.03 & 1,658 & 0.28 & 0.11 & 7,375 & 0.89 & 0.00 \\
Portugal & 26 & 0.80 & 0.03 & 363 & 0.73 & 0.04 & 588 & 0.96 & 0.00 \\
Spain & 137 & 0.68 & 0.09 & 2,240 & 0.68 & 0.03 & 4,047 & 0.93 & 0.00 \\
Denmark & 47 & 0.74 & 0.02 & 688 & 0.74 & 0.03 & 1,934 & 0.87 & 0.00 \\
Israel & 27 & 0.90 & 0.00 & 503 & 0.91 & 0.00 & 784 & 0.89 & 0.00 \\
Norway & 31 & 0.48 & 0.00 & 221 & 0.15 & 0.05 & 1,138 & 0.92 & 0.00 \\
Sweden & 124 & 0.46 & 0.03 & 486 & 0.44 & 0.06 & 3,462 & 0.91 & 0.00 \\
Switzerland & 67 & 0.86 & 0.00 & 546 & 0.79 & 0.01 & 3,734 & 0.73 & 0.00 \\
United Kingdom & 331 & 0.57 & 0.02 & 4,388 & 0.69 & 0.02 & 6,573 & 0.74 & 0.01 \\
\multicolumn{10}{l}{\textit{Developed markets: Pacific}} \\
Australia & 279 & 0.83 & 0.00 & 1,706 & 0.69 & 0.03 & 1,721 & 0.75 & 0.00 \\
Hong Kong & 206 & 0.89 & 0.02 & 0 & 1.00 & 1.00 & 6,130 & 0.88 & 0.00 \\
Japan & 1,834 & 0.86 & 0.04 & 13,080 & 0.96 & 0.00 & 12,563 & 0.87 & 0.00 \\
New Zealand & 4 & 0.00 & 0.21 & 70 & 0.33 & 0.03 & 132 & 0.69 & 0.00 \\
Singapore & 160 & 0.64 & 0.01 & 211 & 0.61 & 0.02 & 653 & 0.77 & 0.00 \\
\multicolumn{10}{l}{\textit{Emerging markets}} \\
Brazil & 450 & 0.99 & 0.00 & 1,572 & 0.95 & 0.00 & 3,141 & 0.93 & 0.00 \\
Colombia & 27 & 1.00 & 0.00 & 179 & 0.72 & 0.00 & 500 & 0.99 & 0.00 \\
Czech Republic & 29 & 0.06 & 0.00 & 107 & 0.74 & 0.00 & 428 & 0.99 & 0.00 \\
Hungary & 7 & 0.94 & 0.00 & 105 & 0.74 & 0.01 & 539 & 0.98 & 0.00 \\
Malaysia & 17 & 0.86 & 0.04 & 338 & 0.84 & 0.01 & 437 & 0.90 & 0.00 \\
Mexico & 91 & 0.97 & 0.02 & 545 & 0.61 & 0.01 & 1,614 & 0.94 & 0.00 \\
Philippines & 9 & 0.85 & 0.05 & 0 & 1.00 & 1.00 & 273 & 0.88 & 0.00 \\
Poland & 42 & 0.98 & 0.00 & 298 & 0.74 & 0.01 & 561 & 0.96 & 0.00 \\
Russia & 9 & 0.96 & 0.00 & 313 & 0.88 & 0.00 & 695 & 0.86 & 0.00 \\
South Africa & 47 & 0.99 & 0.00 & 285 & 0.85 & 0.00 & 1,052 & 0.90 & 0.00 \\
South Korea & 307 & 0.96 & 0.00 & 2,073 & 0.93 & 0.01 & 3,167 & 0.84 & 0.00 \\
Thailand & 82 & 0.98 & 0.01 & 386 & 0.94 & 0.00 & 543 & 0.88 & 0.00 \\
\end{longtable}
\end{landscape}

\subsection{Table 2: Top Ten Investors by Asset Class (2020, API)}

\begin{table}[ht]
\centering
\caption{Top ten investors by asset class, 2020 (API replication). Rankings based on IMF CPIS bilateral holdings from the PIP API. Compare with Table~\ref{tab:t2_orig}: the top investors are broadly similar, validating the free-data pipeline.}
\label{tab:t2_2020}
\small
\begin{tabular}{l r l r l r}
\hline\hline
\multicolumn{2}{c}{\textit{Short-term debt}} & \multicolumn{2}{c}{\textit{Long-term debt}} & \multicolumn{2}{c}{\textit{Equity}} \\
\cmidrule(lr){1-2}\cmidrule(lr){3-4}\cmidrule(lr){5-6}
Country & Bn USD & Country & Bn USD & Country & Bn USD \\
\hline
United States & 5,962 & United States & 38,706 & United States & 81,452 \\
Japan & 1,624 & Japan & 15,539 & Japan & 12,025 \\
France & 881 & United Kingdom & 6,071 & China & 11,892 \\
United Kingdom & 463 & France & 5,568 & France & 11,288 \\
Brazil & 448 & Germany & 5,446 & Germany & 8,059 \\
Canada & 417 & Italy & 3,654 & Netherlands & 7,546 \\
Hong Kong & 394 & Canada & 3,087 & Canada & 7,300 \\
South Korea & 309 & South Korea & 2,350 & United Kingdom & 6,337 \\
Germany & 301 & Netherlands & 2,107 & Hong Kong & 5,838 \\
Australia & 292 & Spain & 1,879 & Spain & 3,911 \\
\hline\hline
\end{tabular}
\end{table}

\newpage

%% ============================================================
%%  SECTION 7 — SUMMARY STATISTICS
%% ============================================================
\section{Summary Statistics}
\label{sec:summary}

This section presents supplementary summary statistics and charts that give
the reader a broader understanding of the underlying data used in the
replication.

\subsection{Cross-Sectional Distribution (Latest Year)}

Table~\ref{tab:summary_stats} summarises the cross-sectional distribution
of outstanding amounts across the 37-country sample for the most recent year.
Long-term debt is the largest asset class by aggregate total; equity shows
the widest cross-country dispersion, driven by the outsized size of the
US and Japanese stock markets.

\input{../_output/summary_stats_table.tex}

\subsection{Time-Series Totals by Asset Class}

\begin{figure}[ht]
\centering
\includegraphics[width=0.92\textwidth]{../_output/summary_stats_chart.png}
\caption{\textbf{Total outstanding financial assets by asset class, 37-country sample.} Short-term debt (blue), long-term debt (green), and equity (pink) in USD trillions from the earliest API vintage through 2024. Long-term debt dominates throughout; equity surged in the 2010s before the post-2021 correction. Sources: OECD T720, BIS WS\_NA\_SEC\_DSS, World Bank WDI.}
\label{fig:ss_ts}
\end{figure}

\subsection{Country-Level Outstanding Amounts}

\begin{figure}[ht]
\centering
\includegraphics[width=0.92\textwidth]{../_output/summary_stats_country_bar.png}
\caption{\textbf{Outstanding financial assets by country, 2024.} Each bar shows a country's total outstanding (USD trillions) split by short-term debt (blue), long-term debt (green), and equity (pink), ranked by total. The United States dwarfs all other economies; equity coverage gaps for several emerging markets reflect reliance on OECD F5 proxy data.}
\label{fig:ss_bar}
\end{figure}

\subsection{Foreign Ownership Shares}

\begin{figure}[ht]
\centering
\includegraphics[width=0.92\textwidth]{../_output/summary_stats_foreign_share.png}
\caption{\textbf{Foreign ownership share by country and asset class, 2024.} Each dot shows the fraction of outstanding assets held by foreign investors per IMF CPIS. Small open economies (Belgium, Netherlands, Singapore) have the highest foreign shares; the United States and Japan have low foreign shares despite their large market size, illustrating the home-bias pattern the paper's demand system models.}
\label{fig:ss_fs}
\end{figure}

\end{document}
